\chapter*{Absolute Convergence, Ratio Test, Root Test, and Alternating Series}
\addcontentsline{toc}{chapter}{Absolute Convergence, Ratio Test, Root Test, and Alternating Series}

\section*{Absolute Convergence}

\begin{definition}
A series $\sum a_n$ is \textbf{absolutely convergent} if
$$\sum |a_n|$$
converges.
\end{definition}

\begin{theorem}
If a series converges absolutely, then it converges.
\end{theorem}

If $\sum a_n$ converges but $\sum |a_n|$ diverges, then the series is called \textbf{conditionally convergent}.

\section*{The Ratio Test}

The ratio test is useful when factorials or exponentials appear.

\begin{theorem}[Ratio Test]
Let $\sum a_n$ be a series with $a_n\neq 0$ eventually, and define
$$L = \lim_{n\to\infty}\left|\frac{a_{n+1}}{a_n}\right|.$$ 
Then:
\begin{itemize}
    \item if $L<1$, the series converges absolutely,
    \item if $L>1$ or $L=\infty$, the series diverges,
    \item if $L=1$, the test is inconclusive.
\end{itemize}
\end{theorem}

Typical examples include
$$\sum \frac{n!}{3^n}, \qquad \sum \frac{x^n}{n!}, \qquad \sum \frac{2^n}{n^3}.$$ 

\section*{The Root Test}

The root test is often cleaner when the $n$th power is built directly into the terms.

\begin{theorem}[Root Test]
Let
$$L = \lim_{n\to\infty}\sqrt[n]{|a_n|}.$$ 
Then:
\begin{itemize}
    \item if $L<1$, the series converges absolutely,
    \item if $L>1$, the series diverges,
    \item if $L=1$, the test is inconclusive.
\end{itemize}
\end{theorem}

This is particularly convenient for terms like $(\frac{2n}{3n+1})^n$.

\section*{Alternating Series}

An alternating series has the form
$$\sum_{n=1}^{\infty} (-1)^{n-1}b_n \qquad \text{or} \qquad \sum_{n=1}^{\infty} (-1)^n b_n,$$
where $b_n\ge 0$.

\begin{theorem}[Alternating Series Test]
If
\begin{enumerate}[label=\arabic*.]
    \item $b_n\ge 0$,
    \item $b_{n+1}\le b_n$ for all sufficiently large $n$,
    \item $\lim_{n\to\infty} b_n = 0$,
\end{enumerate}
then the alternating series converges.
\end{theorem}

\subsection*{Error estimate}

If $S$ is the exact sum and $S_N$ is the $N$th partial sum, then for an alternating series satisfying the test,
$$|S-S_N| \le b_{N+1}.$$ 
So the first omitted term bounds the error.

\begin{remark}
A convergent alternating series need not converge absolutely. The classic example is the alternating harmonic series
$$\sum_{n=1}^{\infty} (-1)^{n-1}\frac1n,$$
which converges conditionally.
\end{remark}